% ch50_infty_adjunctions.tex — Volume IV Chapter 14

\chapter{$\infty$-Adjunctions}

\begin{overviewbox}
Adjunctions have appeared in three earlier guises:
\begin{itemize}
  \item Chapter~7 — hom-set bijection $\mathcal{D}(LX, Y) \cong \mathcal{C}(X, RY)$.
  \item Chapter~17 — triangle identities; uniqueness via Yoneda; adjoint triples.
  \item Chapter~33 — formal account in 2-categories.
\end{itemize}
This chapter develops the fourth incarnation: the $\infty$-categorical
version, where the hom-bijection is replaced by an equivalence of
mapping spaces.

A pair of $\infty$-functors $L : \mathcal{C} \to \mathcal{D}$ and
$R : \mathcal{D} \to \mathcal{C}$ form an \term{$\infty$-adjunction}
$L \dashv R$ if there is a natural equivalence
\begin{equation*}
  \mathrm{Map}_\mathcal{D}(L X, Y) \;\simeq\; \mathrm{Map}_\mathcal{C}(X, R Y),
  \quad X \in \mathcal{C}, \, Y \in \mathcal{D}.
\end{equation*}
The naturality is in the $\infty$-categorical sense: the equivalences fit
together into a coherent system, witnessing arbitrarily high-dimensional
compatibilities.

This chapter develops three equivalent characterizations: hom-equivalence,
unit/counit data, and ``relative adjunctions'' via Cartesian fibrations.
The latter is the model-independent definition (Riehl–Verity) and the
most flexible for proofs.

Theorems: right adjoints preserve $\infty$-limits, left adjoints preserve
$\infty$-colimits (Chapter~49's RAPL/LAPC). Uniqueness up to contractible
choice: any two adjoints to $L$ are equivalent via a unique-up-to-coherent-homotopy
$\infty$-natural isomorphism. The $\infty$-adjoint functor theorem
(presentable case) is stated in §50.4, with proof deferred to Chapter~52.

Adjunctions are the backbone of practical $\infty$-categorical
constructions: free/forgetful adjunctions, localizations, fibre/cofibre
functors, base-change pairs in derived algebraic geometry. By Chapter~52
(presentable case), the adjoint functor theorem will become as central in
the $\infty$-setting as in the $1$-setting.
\end{overviewbox}


% ═══════════════════════════════════════════════════════════════════════════
\section{Defining $\infty$-Adjunctions}
\label{sec:defining-adjunctions}
% ═══════════════════════════════════════════════════════════════════════════

\subsection{Formalizing}

\begin{definition}[$\infty$-Adjunction: hom-equivalence form]
\label{def:adjunction-hom}
A pair of $\infty$-functors $L : \mathcal{C} \to \mathcal{D}$ and
$R : \mathcal{D} \to \mathcal{C}$ form an \term{$\infty$-adjunction}
$L \dashv R$ if there is a natural equivalence of mapping spaces
\begin{equation*}
  \mathrm{Map}_\mathcal{D}(L X, Y) \;\simeq\; \mathrm{Map}_\mathcal{C}(X, R Y)
\end{equation*}
in the variables $X \in \mathcal{C}$ and $Y \in \mathcal{D}$.

Precisely: the two functors $\mathrm{Map}_\mathcal{D}(L (-), -)$ and
$\mathrm{Map}_\mathcal{C}(-, R(-))$ from $\mathcal{C}^{\mathrm{op}} \times
\mathcal{D}$ to $\mathcal{S}$ are equivalent in
$\mathrm{Fun}(\mathcal{C}^{\mathrm{op}} \times \mathcal{D}, \mathcal{S})$.
\end{definition}

\begin{howtosay}
$L \dashv R$ \sayit{``$L$ is left adjoint to $R$''} or \sayit{``$L$
adjoint $R$''}; $L$ called the \term{left adjoint}, $R$ the
\term{right adjoint}.
\end{howtosay}

\begin{definition}[$\infty$-Adjunction: unit/counit form]
\label{def:adjunction-unit-counit}
Equivalently, an $\infty$-adjunction is a pair $L \dashv R$ together with:
\begin{itemize}
  \item A \term{unit} natural transformation $\eta : \mathrm{id}_\mathcal{C}
        \Rightarrow R L$.
  \item A \term{counit} natural transformation $\epsilon : L R \Rightarrow
        \mathrm{id}_\mathcal{D}$.
\end{itemize}
satisfying the \term{triangle identities up to coherent homotopy}:
$\epsilon L \circ L \eta \simeq \mathrm{id}_L$ and $R \epsilon \circ
\eta R \simeq \mathrm{id}_R$, with all higher coherences supplied as
data.
\end{definition}

\begin{theorem}[Equivalence of the two forms]
\label{thm:adjunction-equiv}
The hom-equivalence form (Definition~\ref{def:adjunction-hom}) and the
unit/counit form (Definition~\ref{def:adjunction-unit-counit}) are
equivalent: there is an equivalence of $\infty$-categories of
``$\infty$-adjunctions presented in form $A$'' and ``form $B$.''
\end{theorem}

\begin{remark}[The triangle identities require coherence data]
In a $1$-category, the triangle identities are equations between
natural transformations. In an $\infty$-category, they are weak
equivalences of $1$-simplices, requiring $2$-simplex witnesses; the
$2$-simplex compatibilities require $3$-simplex witnesses; and so on.
The infinite tower of coherences is exactly what distinguishes
$\infty$-adjunctions from naive ``pseudo-adjunctions.''
\end{remark}


% ═══════════════════════════════════════════════════════════════════════════
\section{Relative Adjunctions via Fibrations}
\label{sec:relative-adjunctions}
% ═══════════════════════════════════════════════════════════════════════════

\subsection{Formalizing}

Riehl–Verity propose a model-independent definition of $\infty$-adjunctions
using only the language of $\infty$-categorical mapping spaces and
fibrations.

\begin{definition}[Adjunction via fibrations]
\label{def:adjunction-fibration}
An $\infty$-adjunction $L \dashv R$ is the data of:
\begin{itemize}
  \item A coCartesian fibration $\mathcal{E} \to \Delta^1$ with fibres
        $\mathcal{E}_0 = \mathcal{C}$ over $0$ and $\mathcal{E}_1 = \mathcal{D}$
        over $1$.
  \item The transport functor $\mathcal{C} \to \mathcal{D}$ is $L$.
\end{itemize}
The fibration is simultaneously a Cartesian fibration iff $L$ has a
right adjoint $R$. The Cartesian transport (in the reverse direction)
recovers $R$, and the bi-fibration structure encodes the entire
adjunction data (including coherences).
\end{definition}

\begin{remark}[Why fibrations are model-independent]
The fibration definition uses only $\infty$-categorical features:
existence of (co)Cartesian morphisms, fibres, transport. It works in any
model of $\infty$-categories (quasi-cats, simplicial cats, complete Segal
spaces). The hom-equivalence and unit/counit forms are quasi-categorical
specializations.
\end{remark}


% ═══════════════════════════════════════════════════════════════════════════
\section{Preservation Properties}
\label{sec:adjunction-preservation}
% ═══════════════════════════════════════════════════════════════════════════

\subsection{Formalizing}

\begin{theorem}[$\infty$-RAPL and $\infty$-LAPC]
\label{thm:rapl-lapc-infty}
Let $L \dashv R : \mathcal{D} \to \mathcal{C}$ be an $\infty$-adjunction.
\begin{itemize}
  \item $R$ preserves all $\infty$-limits that exist in $\mathcal{D}$.
  \item $L$ preserves all $\infty$-colimits that exist in $\mathcal{C}$.
\end{itemize}
\end{theorem}

\begin{proof}
Proven already in Theorem~49.\ref{thm:RAPL-infty} as Chapter~49's
$\mathrm{Map}$-Yoneda calculation. \qed
\end{proof}

\begin{remark}[The converse is not automatic]
Preservation of $\infty$-limits does \emph{not} imply being a right adjoint
without further hypotheses (e.g., presentability of $\mathcal{D}$). The
adjoint functor theorem, generalized to the $\infty$-setting, gives a
clean criterion in the presentable case (\S\ref{sec:aft-presentable}).
\end{remark}

\subsection{Reorganizing: uniqueness of adjoints}

\begin{theorem}[Uniqueness of adjoints]
\label{thm:uniqueness-adjoints}
If $L_1 \dashv R$ and $L_2 \dashv R$ are two $\infty$-adjunctions sharing
the same right adjoint, then there is an equivalence $L_1 \simeq L_2$
unique up to a contractible space of equivalences.

Dually: given $L \dashv R_1$ and $L \dashv R_2$, $R_1 \simeq R_2$
uniquely up to contractible choice.
\end{theorem}

\begin{proof}[Proof sketch]
Both $L_1, L_2$ co-represent the functor $Y \mapsto \mathrm{Map}_\mathcal{C}(-, RY)$.
By $\infty$-Yoneda, the corepresenting object is unique up to contractible
space of equivalences. \qed
\end{proof}

\begin{remark}[Adjoint triples]
A right adjoint $R$ may itself have a further right adjoint $R'$, giving
an adjoint triple $L \dashv R \dashv R'$. This phenomenon is essentially
the same as in $1$-categories (Chapter~17), but generalizes: in
$\mathcal{C}\mathrm{at}_\infty$ one finds chains of arbitrary length
(``adjoint chains''), e.g., in the context of $6$-functor formalisms or
derived algebraic geometry.
\end{remark}


% ═══════════════════════════════════════════════════════════════════════════
\section{Adjoint Functor Theorem (Preview)}
\label{sec:aft-presentable}
% ═══════════════════════════════════════════════════════════════════════════

\subsection{Formalizing}

\begin{theorem}[$\infty$-Adjoint Functor Theorem: presentable case]
\label{thm:aft-presentable}
Let $\mathcal{C}$ and $\mathcal{D}$ be presentable $\infty$-categories
(Chapter~52). A functor $R : \mathcal{D} \to \mathcal{C}$ admits a left
adjoint iff $R$ preserves all small $\infty$-limits and is
\term{accessible}: there exists a regular cardinal $\kappa$ such that
$R$ preserves $\kappa$-filtered colimits.

Dually, $L : \mathcal{C} \to \mathcal{D}$ admits a right adjoint iff $L$
preserves all small $\infty$-colimits.
\end{theorem}

\begin{proof}[Proof sketch]
The proof is the central content of HTT §5.5; we defer the full apparatus
to Chapter~52. Key ingredients:
\begin{itemize}
  \item Presentable $\infty$-categories admit accessible localizations.
  \item Accessibility plus colimit preservation gives the construction
        of the adjoint via Kan extension along a presentation.
  \item Uniqueness up to contractible choice follows from $\infty$-Yoneda.
\end{itemize}
\qed
\end{proof}

\begin{remark}[Working horse of higher algebra]
Theorem~\ref{thm:aft-presentable} is the engine of much of higher
algebra: existence of left adjoints to forgetful functors (giving free
constructions), right adjoints to restriction functors (giving induction),
all derive from this single theorem applied to presentable categories.
\end{remark}


% ═══════════════════════════════════════════════════════════════════════════
\section{Examples}
\label{sec:adjunction-examples}
% ═══════════════════════════════════════════════════════════════════════════

\subsection{Foundational adjunctions in Volume~IV}

\begin{example_thm}[$\mathrm{ho} \dashv \mathrm{N}$]
\label{ex:ho-N-adj}
The homotopy category functor $\mathrm{ho} : \mathbf{QCat} \to \mathbf{Cat}$
is left adjoint to the nerve $\mathrm{N} : \mathbf{Cat} \to \mathbf{QCat}$
(Theorem~43.\ref{cor:ho-functor}).
\end{example_thm}

\begin{example_thm}[$\lvert{-}\rvert \dashv \mathrm{Sing}$]
\label{ex:realization-sing-adj}
Geometric realization is left adjoint to the singular complex functor
(Theorem~39.\ref{thm:realization-sing-adj}). The adjunction is a
Quillen equivalence (Theorem~42.\ref{thm:real-sing-qeq}); on
$\infty$-category level, it becomes an equivalence
$\mathcal{S} \simeq \mathcal{S}_{\mathrm{Top}}$.
\end{example_thm}

\begin{example_thm}[Yoneda and free cocompletion]
\label{ex:yoneda-adj}
For any cocomplete $\infty$-category $\mathcal{D}$ and small $\mathcal{C}$,
the restriction $y^* : \mathrm{Fun}^{\mathrm{cc}}(\mathcal{P}(\mathcal{C}),
\mathcal{D}) \to \mathrm{Fun}(\mathcal{C}, \mathcal{D})$ is an equivalence
(Theorem~48.\ref{thm:free-cocompletion}); in particular, it has both
adjoints. The unique colimit-preserving extension $\bar F : \mathcal{P}(\mathcal{C})
\to \mathcal{D}$ is the left adjoint of $y^*$.
\end{example_thm}

\subsection{Reorganizing: where this goes}

\begin{remark}[Adjunctions throughout Volume~IV]
$\infty$-adjunctions appear in nearly every remaining chapter:
\begin{itemize}
  \item Chapter~51: every $\infty$-monad arises from an $\infty$-adjunction;
        $\infty$-Kan extensions are universal adjunctions.
  \item Chapter~52: presentable $\infty$-categories are characterized by
        ``accessibility plus AFT-style adjoint existence.''
  \item Chapter~53: stable $\infty$-cats use adjoints to define
        suspension/loop functors.
  \item Chapter~55: $\infty$-operads are themselves built using
        adjunctions between operadic and ordinary $\infty$-categories.
\end{itemize}
The $\infty$-adjunction is the second pillar of $\infty$-category
theory, after the $\infty$-Yoneda lemma.
\end{remark}


% ═══════════════════════════════════════════════════════════════════════════
\section{Exercises}
\label{sec:adjunctions-exercises}
% ═══════════════════════════════════════════════════════════════════════════

\begin{exercisesbox}
\begin{qa}
\qaitem{Define an $\infty$-adjunction via hom-equivalence.}{$L \dashv R$ iff there's a natural equivalence $\mathrm{Map}_\mathcal{D}(LX, Y) \simeq \mathrm{Map}_\mathcal{C}(X, RY)$ as functors $\mathcal{C}^{\mathrm{op}} \times \mathcal{D} \to \mathcal{S}$.}
\qaitem{What's the unit/counit form?}{$\eta : \mathrm{id} \Rightarrow RL$, $\epsilon : LR \Rightarrow \mathrm{id}$, satisfying triangle identities up to coherent homotopy with infinitely many higher coherences.}
\qaitem{What's the fibrational form?}{A coCartesian fibration $\mathcal{E} \to \Delta^1$ with fibres $\mathcal{C}, \mathcal{D}$ that is also a Cartesian fibration. The transport in two directions gives $L$ and $R$.}
\qaitem{Why is the fibrational form model-independent?}{Because it uses only fibration data, available in any model of $\infty$-categories.}
\qaitem{State $\infty$-RAPL/LAPC.}{$L \dashv R \Rightarrow R$ preserves $\infty$-limits and $L$ preserves $\infty$-colimits. Theorem~49.\ref{thm:RAPL-infty}.}
\qaitem{Why is the proof of RAPL one short calculation?}{$\mathrm{Map}_\mathcal{C}(c, R(\lim p)) \simeq \mathrm{Map}_\mathcal{D}(Lc, \lim p) \simeq \lim \mathrm{Map}_\mathcal{D}(Lc, p) \simeq \lim \mathrm{Map}_\mathcal{C}(c, Rp) \simeq \mathrm{Map}_\mathcal{C}(c, \lim Rp)$. Yoneda concludes.}
\qaitem{Does preserving limits imply being a right adjoint?}{Only in special cases (presentability). The full $\infty$-AFT theorem (\S\ref{sec:aft-presentable}) says: right adjoints among accessible functors between presentable $\infty$-cats = limit-preserving functors.}
\qaitem{State uniqueness of adjoints.}{Two left adjoints to the same $R$ are equivalent up to a contractible space of equivalences. Same for right adjoints.}
\qaitem{Why uniqueness?}{Both adjoints corepresent the same functor $\mathrm{Map}_\mathcal{C}(-, RY)$; $\infty$-Yoneda gives the corepresenting object up to contractible choice.}
\qaitem{What's an adjoint triple?}{$L \dashv R \dashv R'$: $R$ has both a left adjoint $L$ and a further right adjoint $R'$.}
\qaitem{Give a foundational $\infty$-adjunction.}{$\mathrm{ho} \dashv \mathrm{N}$: $\mathbf{QCat} \rightleftarrows \mathbf{Cat}$. Going up: $\mathrm{N}$ embeds $1$-cats into $\infty$-cats; going down: $\mathrm{ho}$ truncates to a $1$-cat.}
\qaitem{Another?}{$\lvert{-}\rvert \dashv \mathrm{Sing}$: realization and singular complex. A Quillen equivalence; on $\infty$-cat level an equivalence $\mathcal{S} \simeq \mathcal{S}_{\mathrm{Top}}$.}
\qaitem{And another?}{$y^* \dashv \bar{(-)}$ via Yoneda: restriction is left adjoint to free-cocompletion extension; gives the universal property of $\mathcal{P}(\mathcal{C})$.}
\qaitem{State the $\infty$-AFT in the presentable case.}{For presentable $\mathcal{C}, \mathcal{D}$: a functor admits a left adjoint iff it preserves all $\infty$-limits and is accessible. Dually for right adjoints.}
\qaitem{What does ``accessible'' mean?}{Preserves $\kappa$-filtered colimits for some regular cardinal $\kappa$. Together with limit preservation, characterizes right adjoints among presentable-target functors.}
\qaitem{Where is $\infty$-AFT proven?}{Chapter~52 (presentability) develops the necessary machinery. Lurie HTT §5.5 has the full account.}
\qaitem{Can $L$ be a right adjoint too?}{Yes — this is called \term{ambidexterity}. Examples: limits and colimits commute in $\mathcal{S}$ (filtered + finite), giving the bifibrational structure.}
\qaitem{How does this generalize Vol I Ch 7 hom-set bijection?}{The hom-set bijection $\mathcal{C}(X, RY) \cong \mathcal{D}(LX, Y)$ becomes a weak equivalence of $\infty$-groupoids (Kan complexes) $\mathrm{Map}$. Naturality is coherent.}
\qaitem{How does it generalize Vol II Ch 17 triangle identities?}{Triangle identities $(\epsilon L)(L\eta) = \mathrm{id}_L$ become weak equivalences with infinitely many higher coherences as data.}
\qaitem{How does it generalize Vol III Ch 33 formal version?}{The $2$-categorical Vol III account becomes the $\infty$-categorical one: $\mathcal{C}\mathrm{at}_\infty$ replaces $\mathbf{Cat}$ as the ambient structure.}
\qaitem{Are $\infty$-adjunctions composable?}{Yes — composition of left adjoints is a left adjoint to the composition of right adjoints. The corresponding fibrations compose via Cartesian product over $\Delta^1$.}
\qaitem{When is an adjunction an equivalence?}{When unit and counit are weak equivalences. Equivalent to: both $L$ and $R$ are fully faithful and essentially surjective.}
\qaitem{What's the next chapter?}{Ch 51: $\infty$-Monads and $\infty$-Kan Extensions. Monads are precisely the algebras associated to $\infty$-adjunctions; Kan extensions are universal adjoints. Together they complete the higher-categorical spine.}
\end{qa}
\end{exercisesbox}


% ═══════════════════════════════════════════════════════════════════════════
\section*{Chapter Summary}
\addcontentsline{toc}{section}{Chapter Summary}
% ═══════════════════════════════════════════════════════════════════════════

\begin{summarybox}
\textbf{$\infty$-Adjunction.}\quad $L \dashv R$ iff there's a natural
equivalence $\mathrm{Map}_\mathcal{D}(LX, Y) \simeq \mathrm{Map}_\mathcal{C}(X, RY)$.
Three forms: hom-equivalence, unit/counit-with-coherences, fibration over
$\Delta^1$.

\textbf{$\infty$-RAPL/LAPC.}\quad Right adjoints preserve $\infty$-limits;
left adjoints preserve $\infty$-colimits. One-line Yoneda proof.

\textbf{Uniqueness.}\quad Any two left adjoints to the same $R$ are
equivalent up to a contractible space of equivalences ($\infty$-Yoneda).

\textbf{$\infty$-AFT (presentable case).}\quad Among accessible functors
between presentable $\infty$-categories, ``right adjoint = preserves
$\infty$-limits.'' Engine of higher-algebra constructions.

\textbf{Foundational examples.}\quad $\mathrm{ho} \dashv \mathrm{N}$;
$\lvert{-}\rvert \dashv \mathrm{Sing}$; Yoneda's $y^* \dashv \bar{(-)}$.

\textbf{Role in Volume~IV.}\quad Used in Chapters~51 (monads), 52
(presentability), 53 (stability), 55 (operads). The second pillar of
$\infty$-CT after Yoneda.
\end{summarybox}


% ═══════════════════════════════════════════════════════════════════════════
\section*{Formal Restatement}
\addcontentsline{toc}{section}{Formal Restatement}
% ═══════════════════════════════════════════════════════════════════════════

\begin{formalrestatement}
\textbf{$\infty$-Adjunction.}\quad $L \dashv R$:
\begin{equation*}
  \mathrm{Map}_\mathcal{D}(LX, Y) \;\simeq\; \mathrm{Map}_\mathcal{C}(X, RY),
  \quad \text{natural in } X, Y.
\end{equation*}

\medskip
\textbf{Unit and counit.}\quad $\eta : \mathrm{id}_\mathcal{C} \Rightarrow RL$,
$\epsilon : LR \Rightarrow \mathrm{id}_\mathcal{D}$. Triangle equivalences:
$\epsilon L \circ L\eta \simeq \mathrm{id}_L$, $R\epsilon \circ \eta R \simeq \mathrm{id}_R$,
plus all higher coherences.

\medskip
\textbf{Fibration form.}\quad $\mathcal{E} \to \Delta^1$ co- and Cartesian
fibration; fibres $\mathcal{C}, \mathcal{D}$; transport in two directions
gives $L, R$.

\medskip
\textbf{Preservation.}\quad $L \dashv R \Rightarrow$ $R$ preserves
$\infty$-limits, $L$ preserves $\infty$-colimits.

\medskip
\textbf{Uniqueness.}\quad The space of adjoints to a given functor is
either empty or contractible.

\medskip
\textbf{$\infty$-AFT (presentable).}\quad For $\mathcal{C}, \mathcal{D}$
presentable and $R : \mathcal{D} \to \mathcal{C}$:
\begin{equation*}
  R \text{ has a left adjoint} \iff R \text{ preserves } \infty\text{-limits and is accessible}.
\end{equation*}
Dually, $L$ has a right adjoint $\iff$ $L$ preserves $\infty$-colimits.
\end{formalrestatement}
